\documentclass[9pt]{extarticle}

\usepackage[utf8]{inputenc}
\usepackage[a4paper, margin=1.4cm]{geometry}
\usepackage[colorlinks=true, urlcolor=blue]{hyperref}
\usepackage{enumitem}
\usepackage{changepage}

\setlist[itemize]{label=\textbullet, noitemsep, topsep=0pt}
\setlength{\parindent}{0pt}

\begin{document}

\begin{center}
    {\huge\bfseries Paul Parisot}
    \\[3pt]
    {\Large Software Engineer - Embedded Linux \& C++ Systems}
    \\[5pt]
\end{center}

\par\vspace{-10pt}
\rule{\linewidth}{0.4pt}
\par\vspace{-4pt}

\begin{center}
    Github: \href{https://github.com/paulogarithm}{paulogarithm} 
    \quad | \quad 
    \href{mailto:paul.parisot@epitech.eu}{paul.parisot@epitech.eu}
\end{center}

\par\vspace{-10pt}
\rule{\linewidth}{0.4pt}
\par\vspace{5pt}

% à propos

{\large\bfseries À propos}
\par\vspace{-8pt}
\rule{\linewidth}{0.4pt}
\par\vspace{5pt}

Ingénieur logiciel orienté systèmes et performance, avec de solides bases en C/C++, Go et Python.
\\
Expérience en interfaçage matériel, traitement de flux de données à haut volume, conception de langages et moteurs de jeu.
\\
Contributeur open-source avec plus de 10 000 lignes de code sur des packages critiques.
\\
En préparation de l'agrégation d'informatique, concours national d'enseignement très sélectif, afin d'approfondir mes bases théoriques en algorithmique des graphes, théorie de la complexité et automates. Un an à McGill (mathématiques, économie, management), GPA 3.57/4.0.
\par\vspace{5pt}

% compétences

{\large\bfseries Compétences}
\par\vspace{-8pt}
\rule{\linewidth}{0.4pt}
\par\vspace{5pt}

\hspace*{5pt}{\bfseries Techniques}: C/C++, Golang, Python, C\#, Flutter, JavaScript/TypeScript, Haskell, Lua(u), OCaml, Assembleur
\\
\hspace*{5pt}{\bfseries Autres}: Outils Git, DevOps (Docker, Kubernetes, ...), CI/CD, Environnement Linux

\par\vspace{5pt}

% expérience professionnelle

{\large\bfseries Expérience Professionnelle}
\par\vspace{-8pt}
\rule{\linewidth}{0.4pt}
\par\vspace{5pt}

{\bfseries Gno.land - Ingénieur Logiciel Junior (Open-Source)}
\begin{itemize}[noitemsep, topsep=0pt]
	\item Développement de 6+ packages core en Gnolang (variante de Go), incluant printf et une compatibilité EVM (conversion d'opcodes EVM en instructions Go).
	\item Détection et correction de {\bfseries 6 bugs critiques} causant des crashs runtime ; {\bfseries ~10k lignes de code} contribuées.
	\item Rédaction de tests exhaustifs sous des politiques strictes de CI/CD et de revue de code.
\end{itemize}

\par\vspace{5pt}

{\bfseries Faurecia / Forvia - Ingénieur Démo Logicielle}
\begin{itemize}[noitemsep, topsep=0pt]
	\item Développement d'une {\bfseries application démo avec Unity (C\#)}, présentée au {\bfseries CES 2024}, simulant navigation GPS \& UI.
	\item Conception d'outils interactifs comme un {\bfseries moteur de calcul d'itinéraire GPS} et une interface de recherche de POI.
	\item {\bfseries Interfaçage matériel} : pilotage de LED en Python, et correction d'un code C++ ancien pour le faire communiquer avec les bons drivers.
	\item Livraison de fonctionnalités prêtes pour la production {\bfseries en avance sur le planning}.
\end{itemize}

\par\vspace{5pt}

{\bfseries Epitech / IONIS - Assistant Technique (Astek) \& Formateur (2 ans)}
\begin{itemize}[noitemsep, topsep=0pt]
	\item Enseignement à {\bfseries plus de 400 étudiants} en C, C++, asm et Haskell ; correction et évaluation des rendus.
	\item Création de {\bfseries contenus pédagogiques interactifs} (quiz, cours inversés, tutoriels).
	\item Application de {\bfseries méthodes pédagogiques} (constructivisme, cognitivisme…) pour renforcer l'engagement des étudiants.
	\item Contribution à un {\bfseries taux de réussite historiquement élevé} au sein des promotions coachées par l'ASTEK.
\end{itemize}

\par\vspace{5pt}

% projets personnels

{\large\bfseries Projets Personnels}
\par\vspace{-8pt}
\rule{\linewidth}{0.4pt}
\par\vspace{5pt}

\underline{\bfseries Merops (ex-Sledge) -- Moteur de migration, synchronisation et orchestration multi-cloud}
\par
\begin{adjustwidth}{5pt}{0pt}
    Développement et finalisation d'un outil générique de migration multi-cloud gérant des structures de données complexes et des interactions avec des bases de données. \\
    Architecture et livraison d'une plateforme backend complète en Go, optimisée pour des opérations de flux de données à haut volume, via une couche d'abstraction cloud personnalisée ({\ttfamily provicol}) supportant AWS S3 et GCP.
\end{adjustwidth}
\par
{\bfseries Compétences :} Go, Python, API Docker, SQL, intégrité des données, flux de données, automatisation système

\par\vspace{8pt}

\underline{\bfseries \href{https://github.com/paulogarithm/cclass}{cclass} -- Moteur de mutation mémoire à l'exécution}
\par
\begin{adjustwidth}{5pt}{0pt}
    Conception d'un framework en C pur modifiant dynamiquement des instructions assembleur compilées directement en mémoire. \\
    Utilisation d'appels système {\bfseries \texttt{mmap}} avec des flags de protection précis pour contourner le surcoût du dispatch par vtable. \\
    Application de règles d'alignement mémoire strictes pour garantir une stabilité absolue sur les cibles d'exécution matérielles.
\end{adjustwidth}
\par
{\bfseries Compétences :} C, systèmes Linux, assembleur x86, mmap, sûreté mémoire, débogage bas niveau

\par\vspace{8pt}

\underline{\bfseries \href{https://github.com/paulogarithm/zzzcript}{ZZZcript} -- Langage fonctionnel statiquement typé}
\par
\begin{adjustwidth}{5pt}{0pt}
	Conception d'un {\bfseries lexer, parser et interpréteur AST} from scratch. \\
	Écriture d'un moteur d'évaluation récursif en {\bfseries Go}, inspiré de Haskell et Golang.
\end{adjustwidth}
\par
{\bfseries Compétences :} Go, Haskell, théorie du parsing, conception d'AST, sémantique fonctionnelle, architecture d'interpréteur

\par\vspace{8pt}

\underline{\bfseries \href{https://github.com/paulogarithm/blorox}{BloRox} -- Moteur de jeu 3D modulaire inspiré de Roblox}
\par
\begin{adjustwidth}{5pt}{0pt}
	Conception d'un moteur de jeu en C avec modules de rendu et de logique chargés à l'exécution. \\
	Intégration de scripting Lua et imitation de l'API Roblox pour la compatibilité. \\
	En développement actif, finalisation prévue à l'arrivée des fonctionnalités de reflection de C++26.
\end{adjustwidth}
\par
{\bfseries Compétences :} C, C++, Lua, architecture ECS, systèmes de plugins, liaison à l'exécution, conception de moteur

\par\vspace{8pt}

% formation

{\large\bfseries Formation}
\par\vspace{-8pt}
\rule{\linewidth}{0.4pt}
\par\vspace{5pt}

{\bfseries Epitech, Paris} \hfill 2022 -- 2027 \\
Diplôme d'ingénieur en technologies de l'information (diplôme de niveau Master) \\
Certification Niveau 7 - Expert en ingénierie logicielle

\vspace{6pt}

{\bfseries McGill University, Montréal} \hfill 2025 -- 2026 \\
Certificat en Management, GPA 3.57/4.0 \\
Cours suivis : Mathématiques, Managerial Economics, Comptabilité, Administration des affaires, Management organisationnel

\end{document}
